\subsection{Algorithm Interface}
\label{sec:arch_algorithm}

The algorithm interface separates the assembled system structure from the execution strategy.
The \texttt{System} owns the components, connections, and structural context.
An algorithm owns the procedure that advances this system in time.
This separation allows the same system definition to be executed with different coordination schemes.

Each algorithm advances the coupled system in macro steps of size~$\Delta t$.
During a macro step, it coordinates input propagation, component stepping, output collection, and any additional treatment required by the system structure.
Explicit algorithms evaluate components according to the available execution order.
Algebraic-loop algorithms add coupled treatment for instantaneous dependencies.
Hybrid algorithms add event detection, event localization, and event dispatch.

All algorithms use the same component and system interfaces.
The user can therefore change the execution strategy without changing the component implementations or the system definition.
The concrete algorithms and their numerical procedures are described in Sections~\ref{sec:impl_algorithms} through~\ref{sec:impl_hybrid}.